\renewcommand{\thesection}{S\arabic{section}}
\setcounter{section}{2}

\section{Frequency-constrained reconstruction and feature analysis}
\label{sec:Appendix_S3}

\subsection{Data and downstream target}
\label{sec:appendix_downstream_target}

We used daily confirmed COVID-19 case incidence from the Johns Hopkins
University Center for Systems Science and Engineering data stream, accessed
through the Delphi Epidata API. The reconstruction analysis was performed on
the national US series from April 1, 2020, through June 30, 2021. For
visualization and calculation of the epidemiological feature summaries, we
restricted the displayed interval to September 1, 2020, through March 31,
2021.

Let
\[
y=(y_1,\ldots,y_T)^\top
\]
denote the observed daily case-incidence series. Because the raw observations
contain day-of-week effects and other short-timescale reporting variation, we
first estimated a smooth downstream mean
\[
\widehat{\mu}
=
(\widehat{\mu}_1,\ldots,\widehat{\mu}_T)^\top
\]
using second-order trend filtering. Specifically, for each target effective
degrees of freedom in
\[
\mathcal D
=
\{10,15,20,30,40,60,80,100,120\},
\]
we fitted a second-order trend-filtering model using the \texttt{genlasso}
implementation in R.

For each fitted mean, we evaluated Gaussian, Poisson, and
negative-binomial observation models. The combination of trend-filtering
degrees of freedom and observation model was selected by the Akaike
information criterion. The fitted mean from the selected combination was used
as the downstream reconstruction target \(\widehat{\mu}\).

Let \(v_t\) denote the fitted conditional variance under the selected
observation model. Thus,
\[
v_t=
\begin{cases}
\widehat{\sigma}^{\,2},
&\text{Gaussian},\\[4pt]
\widehat{\mu}_t,
&\text{Poisson},\\[4pt]
\widehat{\mu}_t+\widehat{\mu}_t^{\,2}/\widehat{\kappa},
&\text{negative binomial},
\end{cases}
\]
where \(\widehat{\kappa}\) is the fitted negative-binomial dispersion
parameter. For the frequency-constrained illustration, we summarized this
possibly time-varying variance by the effective variance
\[
\sigma_{\mathrm{eff}}^2
=
\frac{1}{T}\sum_{t=1}^{T}v_t.
\]
This scalar variance provides a common quadratic scale for constructing and
visualizing the locally non-identifiable perturbation set. It is not intended
to replace the observation-model-specific bounds derived in
Appendix~\ref{sec:Appendix_S1}.


\subsection{Delay kernel and discrete observation operator}
\label{sec:appendix_delay_operator}

The reconstruction used an empirical symptom-onset-to-case-report delay
distribution for COVID-19 reported by Zhang et al. The fitted continuous
lognormal distribution was discretized into daily probabilities,
\[
g_k
=
\Pr(k\leq \tau<k+1),
\qquad k=0,\ldots,60,
\]
and renormalized so that
\[
\sum_{k=0}^{60}g_k=1.
\]

For an upstream incidence vector
\[
U=(U_1,\ldots,U_T)^\top,
\]
the corresponding downstream mean is
\[
KU,
\]
where \(K\in\mathbb R^{T\times T}\) is the lower-triangular causal convolution
matrix with entries
\[
K_{ts}
=
\begin{cases}
g_{t-s},&0\leq t-s\leq 60,\\
0,&\text{otherwise}.
\end{cases}
\]
Equivalently,
\[
(KU)_t
=
\sum_{k=0}^{\min(60,t-1)}g_kU_{t-k}.
\]
The finite-series calculation therefore uses causal linear convolution rather
than circular convolution.

\subsection{Frequency-constrained upstream reconstructions}
\label{sec:appendix_cutoff_reconstruction}

To isolate the effect of assumptions about admissible temporal scales, we
represented the upstream trajectory in a real Fourier basis. Let
\[
\Phi_{f_c}
=
\begin{bmatrix}
\phi_1 & \cdots & \phi_{q(f_c)}
\end{bmatrix}
\in\mathbb R^{T\times q(f_c)}
\]
contain the normalized constant, sine, and cosine basis vectors whose
frequencies satisfy
\[
0\leq f\leq f_c.
\]
The basis vectors were normalized to have unit Euclidean norm. Apart from the
constant component, each non-Nyquist Fourier frequency contributes one sine
and one cosine basis vector.

An upstream trajectory in the cutoff-specific class is written as
\[
U=\Phi_{f_c}\theta,
\]
where \(\theta\in\mathbb R^{q(f_c)}\) is the vector of Fourier coefficients.
For each cutoff \(f_c\), the reference reconstruction was obtained by solving
\[
\widehat{\theta}_{f_c}
=
\arg\min_{\theta}
\left\{
\left\|
K\Phi_{f_c}\theta-\widehat{\mu}
\right\|_2^2
+
\lambda_\theta
\sum_{j=2}^{q(f_c)}\theta_j^2
\right\},
\]
subject to
\[
\Phi_{f_c}\theta\geq 0.
\]
The first coefficient, corresponding to the constant component, was excluded
from the ridge penalty. We used the small numerical ridge parameter
\[
\lambda_\theta=10^{-6},
\]
which stabilizes the optimization without materially determining the shape of
the solution.

The resulting cutoff-specific reconstruction and its delayed projection are
\[
U^\star_{f_c}
=
\Phi_{f_c}\widehat{\theta}_{f_c},
\qquad
D^\star_{f_c}
=
KU^\star_{f_c}.
\]

We considered
\[
f_c\in
\left\{
\frac{1}{35},
\frac{1}{14},
\frac{1}{7}
\right\}
\quad\text{cycles/day}.
\]
These cutoffs admit variation with periods no shorter than 35, 14, and 7 days,
respectively. The three trajectories are therefore not estimates from three
competing epidemiological models. They form a controlled sequence of
reconstruction classes that progressively permits additional short-timescale
structure.

\subsection{Frequency-constrained non-identifiable perturbation bands}
\label{sec:appendix_perturbation_bands}

For each cutoff-specific reference trajectory, we considered perturbations
within the same Fourier subspace,
\[
\delta U
=
\Phi_{f_c}a,
\]
where \(a\in\mathbb R^{q(f_c)}\) denotes the perturbation coefficients. Under
the effective homoscedastic quadratic approximation used for this
visualization, the downstream separability of a perturbation is
\[
J(\delta U)
=
\frac{\|K\delta U\|_2^2}
{\sigma_{\mathrm{eff}}^2}.
\]

To obtain an analytically interpretable frequency-domain approximation to this
quantity, let
\[
G_T(f)
=
\sum_{k=0}^{60}g_k
\exp(-2\pi i f k)
\]
denote the discrete Fourier response of the delay kernel evaluated on the
length-\(T\) Fourier grid. For a normalized sine or cosine basis component at
frequency \(f_j\), define
\[
\lambda_j
=
\frac{|G_T(f_j)|^2}
{\sigma_{\mathrm{eff}}^2}.
\]
The separability of a perturbation is then approximated by the diagonal
spectral quadratic form
\[
J(\delta U)
\approx
a^\top\Lambda_{f_c}a,
\qquad
\Lambda_{f_c}
=
\operatorname{diag}(\lambda_1,\ldots,\lambda_{q(f_c)}).
\]

For a target total testing error \(\alpha=0.20\), the Gaussian
non-identifiability threshold derived in
Appendix~\ref{sec:Appendix_S1} is
\[
\tau_\alpha
=
4(1-\alpha)^2
=
2.56.
\]
We therefore define the cutoff-specific perturbation ellipsoid
\[
\mathcal A_{f_c}
=
\left\{
a:
a^\top\Lambda_{f_c}a
\leq
\tau_\alpha
\right\}.
\]

At each time \(t\), the largest possible absolute perturbation within this
ellipsoid is
\[
B_{f_c}(t)
=
\sup_{a\in\mathcal A_{f_c}}
\left|
\phi_{f_c}(t)^\top a
\right|,
\]
where
\[
\phi_{f_c}(t)
=
\left(
\phi_1(t),\ldots,\phi_{q(f_c)}(t)
\right)^\top.
\]
By the Cauchy--Schwarz inequality in the
\(\Lambda_{f_c}\)-weighted norm,
\[
B_{f_c}(t)
=
\sqrt{
\tau_\alpha\,
\phi_{f_c}(t)^\top
\Lambda_{f_c}^{-1}
\phi_{f_c}(t)
}.
\]
Equivalently,
\[
B_{f_c}(t)
=
\left[
\tau_\alpha
\sum_{j=1}^{q(f_c)}
\frac{\phi_j(t)^2}{\lambda_j}
\right]^{1/2}.
\]

The band shown in Fig.~\ref{fig:epi-features}A is
\[
\left[
\max\left\{
U^\star_{f_c}(t)-B_{f_c}(t),0
\right\},
\;
U^\star_{f_c}(t)+B_{f_c}(t)
\right].
\]
The lower boundary was truncated at zero to respect the nonnegativity of
incidence.

The band is a pointwise envelope of the perturbation ellipsoid rather than a
simultaneous confidence or credible band. In particular, the perturbation
attaining the upper or lower endpoint may differ across time points. The band
therefore visualizes the local geometry of the non-identifiable upstream set;
it should not be interpreted as the sampling distribution of a specific
reconstruction estimator.


\subsection{Epidemiological feature summaries}
\label{sec:Appendix_S4}

We evaluated three quantities derived from each cutoff-specific upstream
trajectory: epidemic peak timing, epidemic peak magnitude, and the seven-day
log growth rate. All feature summaries were calculated over the displayed
interval from September 1, 2020, through March 31, 2021.

\subsubsection{Peak timing and peak magnitude}

For each cutoff-specific reconstruction, we first formed a 14-day centered
moving average,
\[
\widetilde U_{f_c}(t)
=
\frac{1}{14}
\sum_{s\in\mathcal W_t}
U^\star_{f_c}(s),
\]
where reflection padding was used at the boundaries of the finite interval.
The peak index was defined as
\[
t^\star_{f_c}
=
\arg\max_t
\widetilde U_{f_c}(t).
\]

For visualization, relative peak timing was calculated as
\[
T^{\mathrm{rel}}_{f_c}
=
\frac{t^\star_{f_c}}
{T_{\mathrm{disp}}-1},
\]
where \(T_{\mathrm{disp}}\) is the number of days in the displayed analysis
interval.

The peak magnitude was evaluated from the unsmoothed cutoff-specific
reconstruction at the peak index selected from the smoothed trajectory,
\[
M_{f_c}
=
U^\star_{f_c}(t^\star_{f_c}).
\]
Relative peak magnitude was then defined as
\[
M^{\mathrm{rel}}_{f_c}
=
\frac{M_{f_c}}
{\max_{f_c'}M_{f_c'}}.
\]
The common 14-day smoothing used to locate the epidemic peak reduces the
sensitivity of the peak summaries to isolated daily fluctuations while
preserving the broad epidemic-wave structure.


\subsubsection{Seven-day log growth rate}
\label{sec:appendix_growth_rate}

For each reconstructed upstream trajectory, we calculated the seven-day log
growth rate
\[
r_{7,f_c}(t)
=
\frac{1}{7}
\left[
\log\left\{
U^\star_{f_c}(t+7)+c
\right\}
-
\log\left\{
U^\star_{f_c}(t)+c
\right\}
\right],
\]
where \(c=1\) is a numerical offset used to avoid taking the logarithm of zero.
The resulting growth-rate trajectory was smoothed using a five-day centered
moving average with reflection padding at the boundaries. This final
smoothing was applied identically to all cutoff-specific trajectories.


\subsection{Spectral sensitivity of the seven-day growth rate}
\label{sec:appendix_growth_spectral}

The seven-day growth rate is nonlinear in \(U(t)\), but its local spectral
sensitivity can be seen through a first-order perturbation. Let
\[
U_\epsilon(t)
=
U(t)+\epsilon\,\delta U(t).
\]
Differentiating \(r_7(t)\) with respect to \(\epsilon\) at \(\epsilon=0\)
gives
\[
\left.
\frac{\partial r_7[U_\epsilon](t)}
{\partial\epsilon}
\right|_{\epsilon=0}
=
\frac{1}{7}
\left[
\frac{\delta U(t+7)}
{U(t+7)+c}
-
\frac{\delta U(t)}
{U(t)+c}
\right].
\]

If the denominator varies slowly over a local interval, so that
\[
U(t+7)+c
\approx
U(t)+c
\approx
q(t),
\]
then
\[
\delta r_7(t)
\approx
\frac{
\delta U(t+7)-\delta U(t)
}{
7q(t)
}.
\]
The corresponding temporal-difference operator has frequency response
\[
H_7(f)
=
\frac{
e^{2\pi i 7f}-1
}{7},
\]
with magnitude
\[
|H_7(f)|
=
\frac{
2|\sin(7\pi f)|
}{7}.
\]
For low frequencies,
\[
|H_7(f)|
\approx
2\pi |f|,
\]
so the operator suppresses nearly constant components and increases its
sensitivity as frequency moves away from zero. Relative to summaries of the
overall level or broad epidemic-wave shape, the growth-rate calculation
therefore places greater emphasis on short-timescale differences in the
reconstructed trajectory.

The five-day moving average subsequently applied to the estimated growth rate
attenuates the highest-frequency numerical variation, but it does not remove
the underlying dependence of local growth estimates on upstream components
that are more weakly transmitted through the delay process.


Peak timing and peak magnitude are nonlinear global functionals and do not
possess a fixed linear transfer function analogous to \(H_7(f)\). Their
relative stability in Fig.~\ref{fig:epi-features} is therefore an empirical
property of the admissible reconstruction family considered here, rather than
a universal theorem that peak summaries are always identifiable. The analysis
shows that, for this epidemic series and these reconstruction classes, the
broad peak structure is considerably more stable than the local growth-rate
trajectory.